






We evaluate the performance of our method on the FormalGeo-7K dataset. Our main research questions are as follows:

\noindent\textbf{RQ1:} Does our method lead to a higher rate of correct proofs generated by the LLM?


\noindent\textbf{RQ2 (Ablation):} What is the individual contribution of different components in our method?

%\noindent\textbf{RQ3:} How often the LLM produce correct proofs given it arrives the correct numeric answer?


\noindent\textbf{A Note on Baselines for Comparison.}
Our goal is to measure how symbolic augmentation improves formal proof generation in \emph{existing LLMs}, rather than compete with specialized geometry solvers. Similarly, we evaluate whether analogy retrieval improves proof generation, rather than optimize the retrieval method itself. We evaluate retrieval in three ways: \emph{intrinsic retriever evaluation} via proof-similarity prediction (\S\ref{subsec:similar_problems_retrieval}); \emph{intrinsic comparison to random retrieval} via theorem coverage (Table~\ref{tab:theorems_coverage}); and \emph{extrinsic comparison to random retrieval} via downstream proof accuracy under matched inference budgets (Table~\ref{tab:main_results_tab}). Comparisons with additional retrieval methods are left to future work.
